\section{Anhang}

\begin{figure}
  \centering
  \includegraphics[width=0.8\textwidth]{Na_Zerfallsschema.png}
  \caption{Natrium-22 Zerfallsschema aus \cite{Anleitung}}
\label{fig:Na_Zerfallsschema}
\end{figure}

\begin{figure}
  \centering
  \includegraphics[width=0.8\textwidth]{Ba_Zerfallsschema.png}
  \caption{Barium-133 Zerfallsschema aus \cite{Anleitung}}
\label{fig:Ba_Zerfallsschema}
\end{figure}

\begin{figure}
  \centering
  \includegraphics[width=0.8\textwidth]{Spektren/Spektrum1.png}
  \caption{Na Spektrum ohne angepasste Verstärkung (Rechte Seite)}
\label{fig:Spektrum1}
\end{figure}

\begin{figure}
  \centering
  \includegraphics[width=0.8\textwidth]{Spektren/Spektrum2.png}
  \caption{Na Spektrum mit angepasster Verstärkung (Rechte Seite), Gauss-Kurve an 511keV-Linie gefittet}
\label{fig:Spektrum2}
\end{figure}

\begin{figure}
  \centering
  \includegraphics[width=0.8\textwidth]{Spektren/Spektrum3.png}
  \caption{Ba Spektrum (Rechte Seite)}
\label{fig:Spektrum3}
\end{figure}

\begin{figure}
  \centering
  \includegraphics[width=0.8\textwidth]{Spektren/Spektrum3_a3.png}
  \caption{Niederenergetischere Teil des Ba Spektrums (Rechte Seite), Gauss-Kurven an Spektrallinien gefittet}
\label{fig:Spektrum3_a}
\end{figure}

\begin{figure}
  \centering
  \includegraphics[width=0.8\textwidth]{Spektren/Spektrum3_b.png}
  \caption{Hochenergetischere Teil des Ba Spektrums (Rechte Seite), Gauss-Kurven an Spektrallinien gefittet}
\label{fig:Spektrum3_b}
\end{figure}

\begin{figure}
  \centering
  \includegraphics[width=0.8\textwidth]{Spektren/Spektrum4.png}
  \caption{Durch SCA isolierte 511keV-Linie des Na-Spektrums (Rechte Seite)}
\label{fig:Spektrum4}
\end{figure}

%%%%%%%%%%%%%%%%%%%%%%%%%%%%%%%%%%%%%%%%%%%%%%%%%%%%%%%%%%%%%%%%%%%%%%%%%%%%%%%%%%%%%%%%
%%%%%%%%%%%%%%%%%%%%%%%%%%%%%%%%%%%%%%%%%%%%%%%%%%%%%%%%%%%%%%%%%%%%%%%%%%%%%%%%%%%%%%%%

\begin{figure}
  \centering
  \includegraphics[width=0.8\textwidth]{Spektren/Spektrum5.png}
  \caption{Na Spektrum ohne angepasste Verstärkung (Linke Seite)}
\label{fig:Spektrum5}
\end{figure}

\begin{figure}
  \centering
  \includegraphics[width=0.8\textwidth]{Spektren/Spektrum6.png}
  \caption{Na Spektrum mit angepasster Verstärkung (Linke Seite), Gauss-Kurve an 511keV-Linie gefittet}
\label{fig:Spektrum6}
\end{figure}

\begin{figure}
  \centering
  \includegraphics[width=0.8\textwidth]{Spektren/Spektrum7.png}
  \caption{Ba Spektrum (Linke Seite)}
\label{fig:Spektrum7}
\end{figure}

\begin{figure}
  \centering
  \includegraphics[width=0.8\textwidth]{Spektren/Spektrum7_a.png}
  \caption{Niederenergetischere Teil des Ba Spektrums (Linke Seite), Gauss-Kurven an Spektrallinien gefittet}
\label{fig:Spektrum7_a}
\end{figure}

\begin{figure}
  \centering
  \includegraphics[width=0.8\textwidth]{Spektren/Spektrum7_b.png}
  \caption{Hochenergetischere Teil des Ba Spektrums (Linke Seite), Gauss-Kurven an Spektrallinien gefittet}
\label{fig:Spektrum7_b}
\end{figure}

\begin{figure}
  \centering
  \includegraphics[width=0.8\textwidth]{Spektren/Spektrum8.png}
  \caption{Durch SCA isolierte 511keV-Linie des Na-Spektrums (Linke Seite)}
\label{fig:Spektrum8}
\end{figure}

%%%%%%%%%%%%%%%%%%%%%%%%%%%%%%%%%%%%%%%%%%%%%%%%%%%%%%%%%%%%%%%%%%%%%%%%%%%%%%%%%%%%%%%%
%%%%%%%%%%%%%%%%%%%%%%%%%%%%%%%%%%%%%%%%%%%%%%%%%%%%%%%%%%%%%%%%%%%%%%%%%%%%%%%%%%%%%%%%

\begin{figure}
  \centering
  \includegraphics[width=0.8\textwidth]{Spektren/Spektrum9.png}
  \caption{Ba Spektrum mit CFD-Schwelle auf minimaler Einstellung (Rechte Seite)}
\label{fig:Spektrum9}
\end{figure}

\begin{figure}
  \centering
  \includegraphics[width=0.8\textwidth]{Spektren/Spektrum10.png}
  \caption{Ba Spektrum mit CFD-Schwelle auf minimaler Einstellung (Linke Seite)}
\label{fig:Spektrum10}
\end{figure}

\begin{figure}
  \centering
  \includegraphics[width=0.8\textwidth]{Spektren/Spektrum11.png}
  \caption{Ba Spektrum mit CFD-Schwelle auf erhöhter Einstellung (Linke Seite)}
\label{fig:Spektrum11}
\end{figure}

\begin{figure}
  \centering
  \includegraphics[width=0.8\textwidth]{Spektren/Spektrum12.png}
  \caption{Na Promptkurven zur Zeitkalibration}
\label{fig:Spektrum12}
\end{figure}

\begin{figure}
  \centering
  \includegraphics[width=0.8\textwidth]{Spektren/Spektrum12_a.png}
  \caption{Einzelne Na Promptkurve mit Gaussfit zur Bestimmung der Standartabweichung}
\label{fig:Spektrum12_a}
\end{figure}

\begin{figure}
  \centering
  \includegraphics[width=0.8\textwidth]{Spektren/Spektrum12_b.png}
  \caption{Na Promptkurven mit Gaussfits zur Bestimmung der Mittelwerte zur Zeitkalibration}
\label{fig:Spektrum12_b}
\end{figure}

\begin{figure}
    \centering
    \begin{subfigure}{0.5\textwidth}
        \centering
        \includegraphics[width=1\textwidth]{Spektren/Spektrum13.png}
        \caption{}
        \label{fig:Spektrum13_}
    \end{subfigure}%
    \begin{subfigure}{0.5\textwidth}
        \centering
        \includegraphics[width=1\textwidth]{Spektren/Spektrum13_a.png}
        \caption{}
        \label{fig:Spektrum13_a}
    \end{subfigure}
    \caption{Durch SCA isolierte 81keV-Linie des Ba-Spektrums (Rechte Seite)}
    \label{fig:Spektrum13}
\end{figure}

\begin{figure}
    \centering
    \begin{subfigure}{0.5\textwidth}
        \centering
        \includegraphics[width=1\textwidth]{Spektren/Spektrum14.png}
        \caption{}
        \label{fig:Spektrum14_}
    \end{subfigure}%
    \begin{subfigure}{0.5\textwidth}
        \centering
        \includegraphics[width=1\textwidth]{Spektren/Spektrum14_a.png}
        \caption{}
        \label{fig:Spektrum14_a}
    \end{subfigure}
    \caption{Durch SCA isolierte 356keV-Linie des Ba-Spektrums (Linke Seite)}
    \label{fig:Spektrum14}
\end{figure}

\begin{figure}
  \centering
  \includegraphics[width=0.8\textwidth]{Lebensdauermessung.png}
  \caption{Lebensdauerkurve des betrachteten $^{133}$Cs-Zustandes in Einheiten der MCA-Kanäle $b$}
\label{fig:Lebensdauermessung}
\end{figure}